\def\xxactivite{Cours}

\def\xxauteur{P. Bessonnat -- Xavier Pessoles}
\fichefalse \proftrue \tdfalse \courstrue

\def\xxnumchapitre{Chapitre 10 \vspace{.2cm}}

\def\xxchapitre{Dictionnaires -- Mémo}

\def\xxcompetences{%
\textsl{%
\textbf{Savoirs et compétences :}\\
\begin{itemize}[label=\ding{112},font=\color{bleuxp}] 
\item Dictionnaires.
\end{itemize}
}}

\def\xxfigures{
%\includegraphics[width=\linewidth]{fractale}
%\\
%\textit{Modèle du pilote hydraulique avec pilotage interactif.}
}%figues de la page de garde


\input{\repRel/Style/pagegarde_cours_minitoc}
\setlength{\columnseprule}{.1pt}

\vspace{2cm}
\pagestyle{fancy}
\thispagestyle{plain}

%Thème : Tris.
% 
%Commentaires :
%\begin{itemize}
%\item algorithmique quadratique : tri par insertion, par sélection;
%\item tri par partition-fusion;
%\item tri par comptage.
%\end{itemize}
%\textit{On fait observer différentes caractéristiques (par exemple2 stable ou non, en place ou non, comparatif ou non ...).}
%



\section{Introduction}

Les dictionnaires sont composés d'un nombre fini d'éléments auxquels on peut accéder par une clé qui fait partie de l'élément.\\
Chaque élément est donc une paire : une clé et une valeur (\textit{key} et \textit{value} en anglais).\\
Les dictionnaires ne sont pas ordonnés, on ne peut pas rechercher un élément à partir de sa position (indice) dans le dictionnaire mais seulement à partir de la clé.\\
Comme les listes, ce sont des objets \textit{itérables} car on peut parcourir leurs éléments à l'aide d'une boucle \texttt{for}.


\section{Syntaxe}

\begin{itemize}
\item Un dictionnaire python est une succession de paires d'objets séparées par une virgule, délimité par des accolades \{ et \}.
\item Une paire est composée d'une clé et d'une valeur.% qui sont des objets de type quelconque (\texttt{integer}, \texttt{float}, \texttt{string}, \texttt{list}, \texttt{tuple}, \texttt{boolean}).
\item Un dictionnaire est dit \underline{mutable} c'est-à-dire que l'on peut en modifier (voire en supprimer) un ou plusieurs éléments.
\item Les clés du dictionnaire doivent être non mutable (les types \texttt{list} et \texttt{dict} sont interdits).
\item Les valeurs du dictionnaire peuvent être des objets de type quelconque (\texttt{integer}, \texttt{float}, \texttt{string}, \texttt{list}, \texttt{tuple}, \texttt{boolean}).
\end{itemize}

\begin{center}
\texttt{dico=\{cle1 : valeur1, cle2: valeur2  ,cle3 : valeur3,..., clen : valeurn\}\\}
\end{center}

\begin{exemple}~\\

\begin{itemize}
\item \texttt{velo=\{'guidon': 1, 'roue': 2, 'derailleur': 21, 'frein': 2\}}
\item \texttt{pikachu=\{'pokemon':'souris', 'taille': 0.4, 'poids': 6, 'type':'electrik', 'talent':['statik','paratonnerre']\}}
%\item \texttt{altitudeRhone=\{(600,700):\verb![!'Joncin'\verb!]!,(700,800):\verb![!'Luere','Pin-Bouchain','Pilon'\verb!]!,(800,900):\verb![!'Brosses','Croix-Casard'\verb!]!\}}
\end{itemize}
\end{exemple}

\section{Manipulation des dictionnaires}
Le dictionnaire vide est désigné par \{\}.
\subsection*{Création}
La création d'un dictionnaire se fait en choisissant un nom et par les signes d'accolades
\begin{python}
mon\_dico=\{\}
\end{python}

\subsection*{Taille}
La taille du dictionnaire est donnée par la fonction prédéfinie \texttt{length} : 
\texttt{len(mon\_dico)}.


\subsection*{Ajout d'un élément ou modification d'une valeur}
L'ajout d'un élément \texttt{clé:valeur} ou la modification d'une valeur si la clé correspondante existe ont la même syntaxe.

\texttt{
mon\_dico['classe']='PTSI'}
\textit{Remarque} : La clé n'est pas modifiable (non mutable) alors que la valeur l'est.

\subsection*{Suppression d'un élément}
La suppression d'un élément du dictionnaire est réalisée par la fonction \texttt{del} en précisant la clé de l'élément à supprimer. La paire \texttt{clé:valeur} est alors supprimée.
\begin{python}
del(mon\_dico['classe'])
\end{python}

\subsection*{Lecture d'une valeur}

\begin{python}
pikachu['taille'] \# renvoie la valeur associée à la clé 'taille' du dictionnaire pikachu.
\end{python}
Cette instruction peut être associée à \texttt{print} pour afficher la valeur ou affectée à une variable.

\subsection*{Parcours du dictionnaire}
On peut parcourir un dictionnaire par ses clés, ses valeurs ou ses éléments \texttt{clé:valeur}. Pour parcourir la totalité du dictionnaire, on utilise une boucle bornée \texttt{for}.


\begin{center}
\begin{tabular}{|l|p{7cm}|}\hline Instruction & Effet \\
\hline
\texttt {for cle in velo.keys():} & parcourt les clés du dictionnaire\\
\hline
\texttt {for valeur in velo.values():} & parcourt les valeurs du dictionnaire\\
\hline
\texttt {for cle,valeur in velo.items():} & parcourt les éléments du dictionnaire\\
\hline
\end{tabular}
\end{center}


\subsection*{Liste des clés ou liste des valeurs}
On peut récupérer les différentes clés ou les différentes valeurs du dictionnaire sous forme de liste.

\begin{center}
\begin{tabular}{|l|p{10cm}|}\hline Instruction & Effet \\
\hline
\texttt {LesCles=list(velo.keys())} & \texttt{['guidon','roue','derailleur','frein']} est affecté à \texttt{LesCles}\\
\hline
\texttt {LesValeurs=list(velo.values())} & \texttt{[1, 2, 21, 2]} est affecté à \texttt{LesValeurs}\\
\hline
\end{tabular}
\end{center}


\subsection*{Vérification d'une clé}
Pour vérifier qu'une clé existe ou non dans un dictionnaire, on utilise le terme d'appartenance \texttt{in}.

\begin{center}
\begin{tabular}{|l|p{11cm}|}\hline Instruction & Effet \\
\hline
\texttt{'guidon' in velo.keys()} & renvoie le booléen \texttt{True} si la clé 'guidon' est dans \texttt{velo}, \texttt{False} sinon\\
\hline
\texttt{'guidon' not in velo.keys()} & renvoie le booléen \texttt{False} si la clé 'guidon' est dans \texttt{velo}, \texttt{True} sinon\\
\hline
\end{tabular}
\end{center}

